% critique.tex — f166r, vv. 294–336 : commentaire critique (articles textuels tirés du commentaire de l'auteur ; répartition : repartition.md)
\cc{312–30}{Chant du chœur.}{Pour l'étude métrique de ce chant du chœur, l'on pourra consulter Lætitia Parker et son ouvrage de référence, \italique{The Songs of Aristophanes}\nc{Lætitia P. E. Parker, \italique{The songs of Aristophanes}, Clarendon Press Oxford (1997), pp. 406-409.}. Elle parvient à faire apparaître comme une sorte de refrain scandé $- \cup \cup - \cup \cup - \cup \cup -$\nc{Lætitia désigne ce vers comme un \italique{énoplien} (dactylo-épitrite prolongé) et Prato comme un \italique{alcméen}.} aux vers 315, 320 et 325.\\
Quoique la découverte soit extrêmement intéressante et encore à creuser, j'ai préféré pour ma part reprendre la colométrie de Prato, plus proche du manuscrit de Ravenne.}
\cc{336–7}{ἢ 'πικηρυκεύεται / Εὐριπίδῃ Μήδοις ἐπὶ βλάβῃ τινὶ.}{Sur la substitution d'Euripide aux Mèdes, voir le commentaire de la pièce. Remplacé ? pas tout-à-fait ; car au côté d'Εὐριπίδῃ continuent de subsister étrangement les Μήδοις... détail qui fleure bon son interpolation : le spectateur athénien n'avait pas besoin qu'on lui souligne ainsi le trait comique...
Or nous possédons un papyrus d'Oxyrhynque\nc{\italique{The Oxyrhynchus Papyri}, IX (1912), ed. Arthur S. Hunt. Il s'agit du papyrus n° 1176, fragment 39, colonne XII 8-16. Ce papyrus date du \nieme{II} siècle ap. J.-C. et comporte un fragment d'une \italique{Vie} d'Euripide par Satyros de Callatis. Je me réfère aussi à l'édition ainsi qu'au commentaire de ce fragment par Stefan Schorn, \italique{Satyros aus Kallatis} (2002).} du \nieme{II} siècle après Jésus-Christ\nc{Quand notre manuscrit de Ravenne est beaucoup plus tardif : \nieme{IX} ou \nieme{X} s. ap. J.-C.} qui comporte une citation des vers 335-337 des \italique{Thesmophories}. Malheureusement, le vers qui nous intéresse est lacunaire :\citComm{ευριπιδιω\douteux{ι}[. .]\douteux{υ}\nc{Schorn indique que la présence d'un \douteux{c} est moins probable, et que ce n'est assurément pas un \douteux{ε}.}τ|επιβλαβη\douteux{ι}τι|νι}Hunt avait déjà remarqué en son temps que la lacune était bien insuffisante pour qu'ait pu s'y trouver Μήδοις τ᾽\nc{Ce qui n'a pas empêché les éditeurs depuis de prétendre dans leur apparat critique que ce papyrus confirmait l'addition du τ᾽ de Biset : Εὐριπίδῃ Μήδοις <τ᾽>... sauf que le τ identifié par Hunt ne pouvait en aucun cas se rattacher aux Μήδοις, absents du papyrus.}.}
